\documentclass[aspectratio=169,10pt]{beamer}

% Load shared preamble with theme, colors, and packages
% ═══════════════════════════════════════════════════════════════════════════════
% SHARED PREAMBLE — Foundation Models Course (HIT)
% To be \input by each lecture file AFTER \documentclass
% ═══════════════════════════════════════════════════════════════════════════════

% ─────────────────────────────────────────────────────────────────────────────
% BEAMER THEME & BASIC SETUP
% ─────────────────────────────────────────────────────────────────────────────
\usetheme{Madrid}
\usecolortheme{default}

% Remove navigation symbols
\beamertemplatenavigationsymbolsempty

% ─────────────────────────────────────────────────────────────────────────────
% COLORS — HIT Branding
% ─────────────────────────────────────────────────────────────────────────────
\definecolor{hitblue}{RGB}{0,51,102}        % Primary: HIT Navy
\definecolor{hitgold}{RGB}{192,148,0}       % Accent: HIT Gold
\definecolor{hitlight}{RGB}{230,240,250}    % Light background

% Override Beamer palette
\setbeamercolor{primary}{fg=white,bg=hitblue}
\setbeamercolor{secondary}{fg=hitblue,bg=hitlight}
\setbeamercolor{structure}{fg=hitblue}
\setbeamercolor{alerted text}{fg=hitgold}

% ─────────────────────────────────────────────────────────────────────────────
% PACKAGES
% ─────────────────────────────────────────────────────────────────────────────
\usepackage{amsmath}
\usepackage{amssymb}
\usepackage{mathtools}
\usepackage{booktabs}
\usepackage{graphicx}
\usepackage{hyperref}
\usepackage{tikz}
\usepackage{pgfplots}
\usepackage{tcolorbox}
\usepackage{listings}
\usepackage{xcolor}
\usepackage{algorithm}
\usepackage{algpseudocode}

% ─────────────────────────────────────────────────────────────────────────────
% TikZ LIBRARIES
% ─────────────────────────────────────────────────────────────────────────────
\usetikzlibrary{shapes,arrows,positioning,calc,chains,scopes}
\pgfplotsset{compat=1.17}

% ─────────────────────────────────────────────────────────────────────────────
% CODE LISTINGS
% ─────────────────────────────────────────────────────────────────────────────
\lstset{
  basicstyle=\ttfamily\small,
  breaklines=true,
  commentstyle=\color{gray},
  keywordstyle=\color{hitblue},
  stringstyle=\color{hitgold},
  showstringspaces=false,
  backgroundcolor=\color{hitlight},
  frame=single,
  rulecolor=\color{hitblue}
}

% ─────────────────────────────────────────────────────────────────────────────
% TCOLORBOX STYLES
% ─────────────────────────────────────────────────────────────────────────────
\tcbset{
  hitbox/.style={
    colback=hitlight,
    colframe=hitblue,
    fonttitle=\bfseries,
    boxrule=1pt
  },
  alertbox/.style={
    colback=yellow!10,
    colframe=hitgold,
    fonttitle=\bfseries,
    boxrule=1.5pt
  }
}

% ─────────────────────────────────────────────────────────────────────────────
% CUSTOM COMMANDS
% ─────────────────────────────────────────────────────────────────────────────

% Placeholder for figures
\newcommand{\placeholder}[3]{
  \begin{center}
    \fbox{\begin{minipage}[c][#2][c]{#1}
      \centering\small\color{gray}[Figure: #3]
    \end{minipage}}
  \end{center}
}

% Section divider frame
\newcommand{\sectionframe}[2]{
  \begin{frame}[plain]
    \begin{tikzpicture}[remember picture,overlay]
      \fill[hitblue] (current page.south west) rectangle (current page.north east);
      \node[anchor=center, white, font=\Large\bfseries] at (current page.center) {
        \begin{tabular}{c}
          #1 \\[0.3cm]
          {\small\color{hitgold}#2}
        \end{tabular}
      };
    \end{tikzpicture}
  \end{frame}
}

% ─────────────────────────────────────────────────────────────────────────────
% LOAD CUSTOM MACROS
% ─────────────────────────────────────────────────────────────────────────────
% ═══════════════════════════════════════════════════════════════════════════════
% SHARED MACROS — Mathematical Notation for Foundation Models Course
% ═══════════════════════════════════════════════════════════════════════════════

% ─────────────────────────────────────────────────────────────────────────────
% ACTIVATION & LAYER FUNCTIONS
% ─────────────────────────────────────────────────────────────────────────────
\newcommand{\softmax}{\mathrm{softmax}}
\newcommand{\sigmoid}{\mathrm{sigmoid}}
\newcommand{\relu}{\mathrm{ReLU}}
\newcommand{\gelu}{\mathrm{GELU}}
\newcommand{\LayerNorm}{\mathrm{LayerNorm}}
\newcommand{\FFN}{\mathrm{FFN}}
\newcommand{\MHA}{\mathrm{MHA}}
\newcommand{\Attn}{\mathrm{Attention}}

% ─────────────────────────────────────────────────────────────────────────────
% ATTENTION COMPONENTS
% ─────────────────────────────────────────────────────────────────────────────
\newcommand{\query}{Q}
\newcommand{\key}{K}
\newcommand{\val}{V}
\newcommand{\dmodel}{d_{\text{model}}}
\newcommand{\dk}{d_k}
\newcommand{\nheads}{h}

% Scaled dot-product attention formula
\newcommand{\SDPA}{
  \Attn(\query, \key, \val) = \softmax\left(\frac{\query \key^T}{\sqrt{\dk}}\right) \val
}

\newcommand{\CrossAttn}{\text{CrossAttention}}

% ─────────────────────────────────────────────────────────────────────────────
% PROBABILITY & EXPECTATION
% ─────────────────────────────────────────────────────────────────────────────
\newcommand{\E}{\mathbb{E}}
\newcommand{\Prob}{\mathbb{P}}
\newcommand{\KL}{\mathrm{KL}}
\newcommand{\ELBO}{\mathrm{ELBO}}
\newcommand{\given}{\mid}

% ─────────────────────────────────────────────────────────────────────────────
% NORMS & OPERATIONS
% ─────────────────────────────────────────────────────────────────────────────
\newcommand{\norm}[1]{\left\lVert #1 \right\rVert}
\newcommand{\abs}[1]{\left| #1 \right|}
\newcommand{\inner}[2]{\langle #1, #2 \rangle}
\newcommand{\T}{^{\top}}

% ─────────────────────────────────────────────────────────────────────────────
% VECTORS & MATRICES
% ─────────────────────────────────────────────────────────────────────────────
\newcommand{\bvec}[1]{\mathbf{#1}}
\newcommand{\mat}[1]{\mathbf{#1}}
\newcommand{\iid}{\stackrel{\text{i.i.d.}}{\sim}}
\newcommand{\defeq}{\coloneq}

% ─────────────────────────────────────────────────────────────────────────────
% LANGUAGE MODELING
% ─────────────────────────────────────────────────────────────────────────────
\newcommand{\vocab}{\mathcal{V}}
\newcommand{\context}{\text{context}}
\newcommand{\token}{t}
\newcommand{\seq}{x}
\newcommand{\ptheta}{p_\theta}
\newcommand{\pref}{p_{\text{ref}}}

% ─────────────────────────────────────────────────────────────────────────────
% RLHF & DPO
% ─────────────────────────────────────────────────────────────────────────────
\newcommand{\piref}{\pi_{\text{ref}}}
\newcommand{\pitheta}{\pi_\theta}
\newcommand{\yw}{y_w}
\newcommand{\yl}{y_l}
\newcommand{\DPOloss}{\mathcal{L}_{\mathrm{DPO}}}

% ─────────────────────────────────────────────────────────────────────────────
% RETRIEVAL & RAG
% ─────────────────────────────────────────────────────────────────────────────
\newcommand{\docs}{\mathcal{D}}
\newcommand{\qvec}{q}
\newcommand{\dvec}{d}

% ─────────────────────────────────────────────────────────────────────────────
% AGENTS & REASONING
% ─────────────────────────────────────────────────────────────────────────────
\newcommand{\obs}{o}
\newcommand{\act}{a}
\newcommand{\thought}{t}
\newcommand{\tools}{\mathcal{T}}
\newcommand{\history}{h}
\newcommand{\concat}{\oplus}

% ─────────────────────────────────────────────────────────────────────────────
% SETS & LOGIC
% ─────────────────────────────────────────────────────────────────────────────
\newcommand{\R}{\mathbb{R}}
\newcommand{\N}{\mathbb{N}}
\newcommand{\Z}{\mathbb{Z}}

\endinput


% ─────────────────────────────────────────────────────────────────────────────
% TITLE & AUTHOR (can be overridden in individual lectures)
% ─────────────────────────────────────────────────────────────────────────────
\author[D. Lederman]{Dror Lederman}
\institute{Holon Institute of Technology (HIT) \\ Data Science Department}
\date{\today}

% ─────────────────────────────────────────────────────────────────────────────
% FOOTER & HEADER
% ─────────────────────────────────────────────────────────────────────────────
\setbeamertemplate{footline}{
  \leavevmode%
  \hbox{%
    \begin{beamercolorbox}[wd=0.33\paperwidth,ht=2.25ex,dp=1ex,center]{author in head/foot}%
      \usebeamerfont{author in head/foot}\insertshortauthor
    \end{beamercolorbox}%
    \begin{beamercolorbox}[wd=0.34\paperwidth,ht=2.25ex,dp=1ex,center]{title in head/foot}%
      \usebeamerfont{title in head/foot}\insertshorttitle
    \end{beamercolorbox}%
    \begin{beamercolorbox}[wd=0.33\paperwidth,ht=2.25ex,dp=1ex,right]{frame number}%
      \usebeamerfont{frame number}\insertframenumber{} / \inserttotalframenumber\hspace*{2ex}
    \end{beamercolorbox}%
  }%
  \vskip0pt%
}

\endinput


% ─────────────────────────────────────────────────────────────────────────────
% LECTURE-SPECIFIC METADATA
% ─────────────────────────────────────────────────────────────────────────────
\title{Foundation Models: Paradigm, Capabilities, and Risks}
\subtitle{Lecture 1: From Deep Learning to Foundation Models}
\author[D. Lederman]{Dror Lederman}
\institute{Holon Institute of Technology (HIT)}
\date{Semester 2025}

\begin{document}

% ─────────────────────────────────────────────────────────────────────────────
% TITLE SLIDE
% ─────────────────────────────────────────────────────────────────────────────
\begin{frame}
  \titlepage
\end{frame}

% ─────────────────────────────────────────────────────────────────────────────
% LEARNING OBJECTIVES
% ─────────────────────────────────────────────────────────────────────────────
\begin{frame}{Learning Objectives}
  By the end of this lecture, you will be able to:

  \begin{itemize}
    \item Define \textbf{Foundation Models} and understand their paradigm.
    \item Trace the evolution of \textbf{transfer learning} in deep learning.
    \item Identify and explain \textbf{emergent capabilities} in large models.
    \item Analyze the role of \textbf{scaling laws} in model performance.
    \item Evaluate both \textbf{opportunities and risks} of Foundation Models.
    \item Understand the course structure and grading components.
  \end{itemize}
\end{frame}

% ─────────────────────────────────────────────────────────────────────────────
% COURSE ROADMAP
% ─────────────────────────────────────────────────────────────────────────────
\begin{frame}{Course Roadmap: Foundation Models}
  \begin{center}
    \small
    \begin{tabular}{lll}
      \toprule
      \textbf{Weeks} & \textbf{Module} & \textbf{Topics} \\
      \midrule
      1--3 & \textbf{Foundations} & Paradigm, Scaling Laws, Architectures \\
      4--6 & \textbf{Training \& Adaptation} & Pretraining, Fine-Tuning, Alignment \\
      7--9 & \textbf{Capabilities} & Reasoning, Multimodal, Retrieval \\
      10--12 & \textbf{Reliability} & Evaluation, Safety, Deployment \\
      13 & \textbf{Synthesis} & Open Problems, Projects \\
      \bottomrule
    \end{tabular}
  \end{center}

  \vspace{0.4cm}

  \textbf{Key Papers:} Bommasani 2021, Kaplan 2020, Ouyang 2022, Rafailov 2023, Lewis 2020, OWASP/NIST 2023
\end{frame}

% ─────────────────────────────────────────────────────────────────────────────
% GRADING BREAKDOWN
% ─────────────────────────────────────────────────────────────────────────────
\begin{frame}{Grading Breakdown}
  \begin{center}
    \Large
    \begin{tabular}{lc}
      \toprule
      \textbf{Component} & \textbf{Weight} \\
      \midrule
      Two Technical Assignments & 20\% \\
      Paper Presentation & 15\% \\
      Final Project & 35\% \\
      Exam & 30\% \\
      \bottomrule
      \textbf{Total} & \textbf{100\%} \\
      \bottomrule
    \end{tabular}
  \end{center}
\end{frame}

% ─────────────────────────────────────────────────────────────────────────────
% SECTION: WHAT ARE FOUNDATION MODELS?
% ─────────────────────────────────────────────────────────────────────────────
\sectionframe{What Are Foundation Models?}{Defining the paradigm shift}

\begin{frame}{The Modern AI Landscape}
  \begin{columns}[T]
    \column{0.50\textwidth}
      \textbf{Traditional ML Approach:}
      \begin{itemize}
        \item Task-specific models
        \item Labeled datasets for each task
        \item Fine-tuning from scratch
        \item Limited generalization
      \end{itemize}

    \column{0.48\textwidth}
      \textbf{Foundation Models Approach:}
      \begin{itemize}
        \item One large model, many tasks
        \item Large unlabeled data at scale
        \item Prompt or light fine-tuning
        \item Surprising generalization
      \end{itemize}
  \end{columns}

  \vspace{0.4cm}

  \begin{tcolorbox}[hitbox, title=The Paradigm Shift]
    From task-specific models to \textbf{general-purpose foundation models}.
  \end{tcolorbox}
\end{frame}

\begin{frame}{Definition: Foundation Models}
  \textbf{Bommasani et al.\ (2021)}

  \vspace{0.3cm}

  \begin{tcolorbox}[hitbox, title=Foundation Model]
    A model trained on \textbf{broad data} at \textbf{scale} such that it can be adapted
    (fine-tuned) to a wide range of downstream tasks.
  \end{tcolorbox}

  \vspace{0.4cm}

  \textbf{Key Characteristics:}

  \begin{enumerate}
    \item \textbf{Trained on broad data} — text, images, code, etc., not task-specific.
    \item \textbf{Large scale} — billions of parameters, trillions of tokens.
    \item \textbf{Adaptable} — can be fine-tuned, prompted, or used as-is.
    \item \textbf{General-purpose} — performs well across many downstream tasks.
  \end{enumerate}
\end{frame}

\begin{frame}{Examples of Foundation Models}
  \begin{center}
    \begin{tabular}{lll}
      \toprule
      \textbf{Model} & \textbf{Release} & \textbf{Key Property} \\
      \midrule
      BERT & 2018 & Masked language modeling \\
      GPT-3 & 2020 & Few-shot in-context learning \\
      CLIP & 2021 & Vision-language alignment \\
      ChatGPT & 2022 & Instruction-following at scale \\
      GPT-4 & 2023 & Multimodal, reasoning \\
      Llama 3 & 2024 & Open-source, efficient \\
      \bottomrule
    \end{tabular}
  \end{center}

  \vspace{0.3cm}

  \textit{These models differ in architecture, training data, and alignment, but share the foundation model paradigm.}
\end{frame}

% ─────────────────────────────────────────────────────────────────────────────
% SECTION: TRANSFER LEARNING EVOLUTION
% ─────────────────────────────────────────────────────────────────────────────
\sectionframe{Transfer Learning Evolution}{From ImageNet to Foundation Models}

\begin{frame}{Transfer Learning Timeline}
  \begin{enumerate}
    \item[\textbf{Phase 1}] \textbf{Task-Specific Models} (pre-2012)
      \begin{itemize}
        \item Hand-crafted features, limited reuse.
      \end{itemize}

    \vspace{0.2cm}

    \item[\textbf{Phase 2}] \textbf{Pretrain-Finetune} (2012--2019)
      \begin{itemize}
        \item ImageNet pretraining for vision, fine-tune for tasks.
        \item BERT, ELMo: pretraining on raw text.
      \end{itemize}

    \vspace{0.2cm}

    \item[\textbf{Phase 3}] \textbf{Foundation Models} (2020+)
      \begin{itemize}
        \item Scale to billions of parameters and trillions of tokens.
        \item In-context learning, prompting, few-shot adaptation.
      \end{itemize}
  \end{enumerate}
\end{frame}

\begin{frame}{From ImageNet to Language Foundations}
  \begin{columns}[T]
    \column{0.50\textwidth}
      \textbf{Vision (2012 onwards):}
      \begin{itemize}
        \item AlexNet wins ImageNet.
        \item Pretraining + fine-tuning becomes standard.
        \item Transfer learning is the norm.
      \end{itemize}

    \column{0.48\textwidth}
      \textbf{Language (2018 onwards):}
      \begin{itemize}
        \item BERT: bidirectional pretraining.
        \item GPT: autoregressive pretraining.
        \item Scale up: from 100M to 100B+ parameters.
      \end{itemize}
  \end{columns}

  \vspace{0.4cm}

  \textbf{Lesson:} Transfer learning works because:
  \begin{itemize}
    \item General patterns emerge at scale.
    \item Language and vision share low-level and high-level structure.
  \end{itemize}
\end{frame}

\begin{frame}{The Scaling Hypothesis}
  \textbf{Core Insight:} Larger models perform better.

  \vspace{0.4cm}

  \begin{tcolorbox}[hitbox, title=Scaling Laws]
    Model performance improves predictably with:
    \begin{itemize}
      \item Number of parameters (model size)
      \item Amount of training data
      \item Compute budget
    \end{itemize}
  \end{tcolorbox}

  \vspace{0.3cm}

  \textbf{Key References:}
  \begin{itemize}
    \item Kaplan et al.\ (2020): Scaling Laws for Language Models
    \item Hoffmann et al.\ (2022): Chinchilla Rule (optimal balance)
  \end{itemize}

  \textit{These will be covered in detail in Lecture 2.}
\end{frame}

% ─────────────────────────────────────────────────────────────────────────────
% SECTION: EMERGENT CAPABILITIES
% ─────────────────────────────────────────────────────────────────────────────
\sectionframe{Emergent Capabilities}{Abilities that appear at scale}

\begin{frame}{What Are Emergent Capabilities?}
  \begin{tcolorbox}[alertbox, title=Emergence]
    Capabilities that appear at larger scales but are absent or negligible in smaller models.
  \end{tcolorbox}

  \vspace{0.4cm}

  \textbf{Examples:}
  \begin{itemize}
    \item \textbf{In-context learning}: Models learn from examples in the prompt.
    \item \textbf{Chain-of-thought reasoning}: Breaking problems into steps.
    \item \textbf{Instruction following}: Understanding varied instructions without fine-tuning.
    \item \textbf{Common sense}: Reasoning about the real world.
  \end{itemize}

  \vspace{0.3cm}

  \textit{Key paper:} Wei et al.\ (2022): Emergent Abilities of Large Language Models
\end{frame}

\begin{frame}{In-Context Learning Example}
  \textbf{Prompt:}
  \begin{quote}
    \small
    English to French:\\
    ``Hello'' → ``Bonjour''\\
    ``Good morning'' → ``Bon matin''\\
    ``How are you?'' → ?
  \end{quote}

  \vspace{0.3cm}

  \textbf{Small model (125M):}
  \begin{quote}
    \small
    Output: Confused, random tokens.
  \end{quote}

  \vspace{0.2cm}

  \textbf{Large model (175B, GPT-3):}
  \begin{quote}
    \small
    Output: ``Comment allez-vous?'' (correct!)
  \end{quote}

  \vspace{0.3cm}

  \textit{The model ``learns'' from examples without gradient updates—purely from the context.}
\end{frame}

\begin{frame}{Why Emergence Matters}
  \begin{tcolorbox}[hitbox, title=Paradigm Implications]
    If capabilities emerge unpredictably at scale, we must:
    \begin{enumerate}
      \item Invest in scale to discover new abilities.
      \item Prepare for unexpected behavior.
      \item Develop evaluation methods to detect capabilities.
      \item Design safety measures for unknown risks.
    \end{enumerate}
  \end{tcolorbox}
\end{frame}

% ─────────────────────────────────────────────────────────────────────────────
% SECTION: OPPORTUNITIES
% ─────────────────────────────────────────────────────────────────────────────
\sectionframe{Opportunities}{Why Foundation Models Matter}

\begin{frame}{Opportunities: Knowledge Work}
  \begin{columns}[T]
    \column{0.50\textwidth}
      \textbf{Writing \& Content:}
      \begin{itemize}
        \item Email drafting
        \item Code generation
        \item Documentation
        \item Research assistance
      \end{itemize}

    \column{0.48\textwidth}
      \textbf{Analysis \& Decision:}
      \begin{itemize}
        \item Data summarization
        \item Question answering
        \item Medical diagnosis
        \item Legal research
      \end{itemize}
  \end{columns}
\end{frame}

\begin{frame}{Opportunities: Democratization}
  \begin{tcolorbox}[hitbox, title=Access \& Democratization]
    Foundation models lower barriers to entry:
    \begin{itemize}
      \item Open-source models (Llama, Mistral, Falcon)
      \item API access (ChatGPT, Claude, Gemini)
      \item Fine-tuning tools (LoRA, QLoRA)
      \item Reduced need for ML expertise per task.
    \end{itemize}
  \end{tcolorbox}

  \vspace{0.3cm}

  \textit{A person with coding skills can now build AI applications without ML training.}
\end{frame}

% ─────────────────────────────────────────────────────────────────────────────
% SECTION: RISKS & CHALLENGES
% ─────────────────────────────────────────────────────────────────────────────
\sectionframe{Risks \& Challenges}{The Dark Side of Scale}

\begin{frame}{Risk 1: Bias \& Fairness}
  \textbf{Problem:}
  Training on broad data captures societal biases.

  \vspace{0.3cm}

  \textbf{Examples:}
  \begin{itemize}
    \item Gender bias in word associations.
    \item Racial bias in image generation.
    \item Cultural stereotypes in text generation.
  \end{itemize}

  \vspace{0.3cm}

  \textbf{Scale Effect:} Larger models amplify biases across more use cases.
\end{frame}

\begin{frame}{Risk 2: Hallucinations \& Misinformation}
  \textbf{Problem:}
  Models generate plausible-sounding but false information.

  \vspace{0.3cm}

  \textbf{Examples:}
  \begin{itemize}
    \item Fabricating citations and sources.
    \item Inventing historical facts.
    \item Confidently explaining wrong concepts.
  \end{itemize}

  \vspace{0.3cm}

  \textbf{Challenge:} No easy way to prevent without sacrificing capability.
\end{frame}

\begin{frame}{Risk 3: Safety \& Misuse}
  \textbf{Misuse Scenarios:}
  \begin{itemize}
    \item Automated generation of malware, phishing, propaganda.
    \item Impersonation and deepfakes.
    \item Unauthorized access to sensitive information.
  \end{itemize}

  \vspace{0.3cm}

  \textbf{Alignment Challenge:}
  Models trained to be ``helpful'' can be manipulated by adversarial prompts.

  \vspace{0.3cm}

  \textit{Mitigation requires both technical safeguards and governance.}
\end{frame}

\begin{frame}{Risk 4: Environmental \& Computational Cost}
  \textbf{Training a Large Model:}
  \begin{itemize}
    \item Millions of GPU/TPU hours
    \item Hundreds of thousands of dollars
    \item Significant carbon footprint
  \end{itemize}

  \vspace{0.3cm}

  \textbf{Inference Cost:}
  \begin{itemize}
    \item Each token costs computation.
    \item High-volume services accumulate costs and carbon.
  \end{itemize}

  \vspace{0.3cm}

  \textit{Efficiency and sustainable ML are active research areas.}
\end{frame}

\begin{frame}{Risk 5: Concentration \& Power}
  \textbf{Who Can Build?}
  \begin{itemize}
    \item Frontier models require billions in compute and data.
    \item Accessible only to large labs (OpenAI, Google, Meta, etc.)
    \item Creates power concentration in AI development.
  \end{itemize}

  \vspace{0.3cm}

  \textbf{Implications:}
  \begin{itemize}
    \item A few organizations shape AI direction globally.
    \item Economic moat: hard for startups to compete.
    \item Governance challenges: who sets norms?
  \end{itemize}
\end{frame}

% ─────────────────────────────────────────────────────────────────────────────
% SECTION: THE COURSE AHEAD
% ─────────────────────────────────────────────────────────────────────────────
\sectionframe{What We'll Learn}{A deep dive into foundation models}

\begin{frame}{Topics by Module}
  \begin{columns}[T]
    \column{0.50\textwidth}
      \textbf{Module 1: Foundations (Weeks 1--3)}
      \begin{enumerate}
        \item Foundation Models Paradigm (this lecture)
        \item Scaling Laws \& Optimal Training
        \item Transformer Architectures
      \end{enumerate}

    \column{0.48\textwidth}
      \textbf{Module 2: Training (Weeks 4--7)}
      \begin{enumerate}
        \setcounter{enumi}{3}
        \item Large-Scale Pretraining
        \item Efficient Adaptation (LoRA, etc.)
        \item Alignment (RLHF, DPO)
      \end{enumerate}
  \end{columns}

  \vspace{0.3cm}

  \begin{columns}[T]
    \column{0.50\textwidth}
      \textbf{Module 3: Capabilities (Weeks 8--10)}
      \begin{enumerate}
        \setcounter{enumi}{6}
        \item Reasoning \& Planning
        \item Multimodal Models
        \item Knowledge-Augmented Models
      \end{enumerate}

    \column{0.48\textwidth}
      \textbf{Module 4: Reliability (Weeks 11--13)}
      \begin{enumerate}
        \setcounter{enumi}{9}
        \item Evaluation \& Benchmarks
        \item Safety \& Security
        \item Future Directions
      \end{enumerate}
  \end{columns}
\end{frame}

\begin{frame}{How We'll Learn}
  \textbf{Format:}
  \begin{itemize}
    \item \textbf{Lectures} — concepts, theory, key papers
    \item \textbf{Readings} — original papers (we'll discuss key ones)
    \item \textbf{Assignments} — hands-on with models (Llama, OpenAI API)
    \item \textbf{Project} — design and evaluate your own system
  \end{itemize}

  \vspace{0.3cm}

  \textbf{Prerequisites Assumed:}
  \begin{itemize}
    \item Deep learning fundamentals (CNNs, RNNs, basic Transformers)
    \item NLP basics (tokenization, embeddings)
    \item Python for machine learning
  \end{itemize}

  \vspace{0.3cm}

  \textit{If rusty, come to office hours or review Week 0 materials.}
\end{frame}

% ─────────────────────────────────────────────────────────────────────────────
% DISCUSSION
% ─────────────────────────────────────────────────────────────────────────────
\begin{frame}{Discussion: Foundation Models in 5 Years}
  \begin{enumerate}
    \item Where do you think foundation models will have the \textbf{most impact}?
    \item What \textbf{risks} worry you most?
    \item Should governments \textbf{regulate} foundation models?
    \item How will \textbf{safety} and \textbf{capability} balance out?
  \end{enumerate}

  \vspace{0.3cm}

  \textit{No right answers---think critically about the implications.}
\end{frame}

% ─────────────────────────────────────────────────────────────────────────────
% KEY TAKEAWAYS
% ─────────────────────────────────────────────────────────────────────────────
\begin{frame}{Key Takeaways}
  \begin{tcolorbox}[hitbox, title=Summary]
    \begin{enumerate}
      \item \textbf{Foundation Models} are large, general-purpose models trained on broad data.
      \item They represent a paradigm shift from task-specific to general-purpose AI.
      \item \textbf{Emergent capabilities} appear unpredictably at scale.
      \item \textbf{Opportunities} are immense: democratization, new capabilities.
      \item \textbf{Risks} are real: bias, hallucination, misuse, environmental impact, concentration.
      \item This course will cover architecture, training, adaptation, and evaluation.
    \end{enumerate}
  \end{tcolorbox}
\end{frame}

% ─────────────────────────────────────────────────────────────────────────────
% NEXT LECTURE
% ─────────────────────────────────────────────────────────────────────────────
\begin{frame}{Next Lecture}
  \Large
  \textbf{Lecture 2: Scaling Laws \& Optimal Training}

  \vspace{0.5cm}

  Topics:
  \begin{itemize}
    \item Power laws of model performance
    \item Compute-optimal training (Chinchilla rule)
    \item Balancing model size, data, and compute
    \item Implications for your own projects
  \end{itemize}

  \vspace{0.5cm}

  \textbf{Reading:}
  \begin{itemize}
    \item Kaplan et al.\ (2020): Scaling Laws for Language Models
    \item Hoffmann et al.\ (2022): Training Compute-Optimal LLMs
  \end{itemize}
\end{frame}

% ─────────────────────────────────────────────────────────────────────────────
% CLOSING
% ─────────────────────────────────────────────────────────────────────────────
\begin{frame}
  \centering
  \Huge
  Questions?

  \vspace{0.5cm}

  \normalsize
  \textit{dror.lederman@constrol.net}
\end{frame}

\end{document}
